\documentclass{article}
\usepackage[utf8]{inputenc}
\usepackage[margin=2cm]{geometry}
\usepackage{amsmath}
\usepackage{needspace}
\usepackage{xcolor}
\usepackage{sectsty}
\usepackage{helvet}
\renewcommand{\familydefault}{\sfdefault}
\sectionfont{\normalfont\fontsize{16}{20}\selectfont}
\subsectionfont{\normalfont\bfseries\fontsize{13}{16}\selectfont}
\setlength{\parindent}{0pt}
\pagestyle{empty}

% ============================================================
% MOLDE: CONTENIDO EXTRA / AMPLIADO POR TEMA
% ------------------------------------------------------------
% Cada seccion es un titulo fijo con su explicacion debajo.
% Borra las secciones que no apliquen a un tema en particular.
% Repite este archivo por cada tema (un extra por cada teoria).
% ============================================================

\begin{document}
\large

\section*{[NOMBRE DEL TEMA]}

\needspace{4\baselineskip}
\subsection*{Idea intuitiva}
[Explicacion]

\needspace{4\baselineskip}
\subsection*{Procedimiento paso a paso (algoritmo)}
[Explicacion]

% Regla para "Errores frecuentes": entre 5 y 10 ejemplos, y solo
% errores comunes/tipicos que comete un estudiante real (no casos
% raros ni de borde). Si no hay al menos 5 errores comunes genuinos
% para este tema, poner los que haya en vez de inventar hasta llegar
% al minimo.
\needspace{4\baselineskip}
\subsection*{Errores frecuentes}
[Explicacion]

\needspace{4\baselineskip}
\subsection*{Cuando usar este metodo}
[Explicacion]

\needspace{4\baselineskip}
\subsection*{Cuando no usar este metodo}
[Explicacion]

\needspace{4\baselineskip}
\subsection*{Casos especiales}
[Explicacion]

\needspace{4\baselineskip}
\subsection*{Trucos o atajos}
[Explicacion]

\needspace{4\baselineskip}
\subsection*{Advertencias importantes}
[Explicacion]

\needspace{4\baselineskip}
\subsection*{Comprobacion del resultado}
[Explicacion]

\needspace{4\baselineskip}
\subsection*{Relacion con otros temas}
[Explicacion]

\needspace{4\baselineskip}
\subsection*{Requisitos previos}
[Explicacion]

\needspace{4\baselineskip}
\subsection*{Que viene despues de este tema}
[Explicacion]

\needspace{4\baselineskip}
\subsection*{Preguntas frecuentes}
[Explicacion]

\needspace{4\baselineskip}
\subsection*{Curiosidades}
[Explicacion]

\needspace{4\baselineskip}
\subsection*{Notacion alternativa}
[Explicacion]

\end{document}